\documentclass{article}
%  ╔══════════════════════════════════════════════════════════╗
%  ║                         Title                            ║
%  ╚══════════════════════════════════════════════════════════╝
%NOTE: I want \title,\author to be near top, and this \renewcommand{\maketitle} fails if these are above it.
\usepackage[utf8]{inputenc}         % Used to compile any UTF-8 Character (and supports Unicode)
\usepackage{titling}

\renewcommand{\maketitle}{
    \quad
    \vspace{1em}
  \begin{center}
      \LARGE\bfseries \thetitle \par
      \vspace{0.6em}
      \large
  \end{center}
    \vspace{1em}
}

\title{Galois Theory - Intuition Behind Field Extensions}
\author{Nathanael Chwojko-Srawley}
\date{\today}

%  ╔══════════════════════════════════════════════════════════╗
%  ║                         packages                         ║
%  ╚══════════════════════════════════════════════════════════╝

\usepackage{extarrows}         % Used to compile any UTF-8 Character (and supports Unicode)
\usepackage{stackengine}
\usepackage{colortbl}
\usepackage{scalerel}
\usepackage{amsmath, amssymb} % Used to write better mathematical expressions (ex. \mathbb{R})
\usepackage{stackrel}         % for left-right arrows, staking symbol above and bellow
\usepackage{mathrsfs}
\usepackage{bbm}                     % for mathbbm{1}
\usepackage{euscript}
\usepackage{commutative-diagrams}
\usepackage[font=itshape]{quoting}
\usepackage{mathtools}
  \usepackage{enumitem}
\usepackage{amsthm}           % Used to create theorem enviroments.
\usepackage{graphicx}         % Let's you use images in LaTeX; ex the \includegraphics command.
\usepackage{hhline}
%\usepackage{luatexja}
%\usepackage{fontspec}
\usepackage{dsfont}
\usepackage{wasysym}          % emoji's!
\usepackage{float}            % package with ``H'' for figures and tables, ex. Images, tables, etc
%\usepackage{subcaption}      % More options with sub-captions.
\usepackage{setspace}         % More flexibility on spacing in document.
\usepackage{hyperref}         % let's me put hyperlinks
%\usepackage{tikz}             % for commutative diagrams https://tex.stackexchange.com/questions/419221/how-to-do-the-pushout-with-universal-property
\usepackage{tikz-cd}          % for commutative diagrams https://tex.stackexchange.com/questions/419221/how-to-do-the-pushout-with-universal-property
%\usepackage{quiver}           % for better commutative diagrams
\usepackage{bookmark}         % creating heading shortcut bar on the left as in word
\usepackage{longtable}        % create tables that extend past one page
%\usepackage{siunitx}         % required for alignment
%\usepackage{multirow}        % for multiple rows
\usepackage{microtype}        % makes nice text. Fixes some under-/over- filled box problems.
%\usepackage{fancyvrb}        % Let's you write code (different style than listings).
\usepackage[most]{tcolorbox} 	      % Makes nice boxes for thms, cor, prop, or anything else.
%\usepackage{enumitem}        % More flexibility on enumeration (don't know much).
\usepackage{xcolor}           % Colour text.
\usepackage{wrapfig}          % To wrap text around figure
\usepackage{listings}         % Create nice colour code blocks (customized in preamble)
\usepackage{thmtools}         % Customize theorem environment (in preamble).
%\usepackage{marvosym}        % Provide ``basic'' symbols (hazard, religious, smiley face etc.)
\usepackage{array}            % \begin{table}{>{\em}c c} -> 1st column italic (bfseries for bold)
%\usepackage[english]{babel}  % If I want to blindtext in english.
\usepackage{blindtext}
\usepackage[a4paper, total={6in, 8.5in}]{geometry}
%\usepackage[symbol]{footmisc} % For daggers and asterisk for footnotes. 1: *, 2: dagger, etc.
\usepackage{fancyhdr}
\usepackage{titlesec}
%\usepackage{pst-plot}          %To make nice plots: https://tex.stackexchange.com/questions/47594/how-can-i-make-a-graph-of-a-function#47596
% \usepackage[answerdelayed]{exercise}
\usepackage{etoolbox}          % for Dynkin diagrams
\usepackage{dynkin-diagrams}   % for Dynkin diagrams

% ╔═════════════════════════════════════════════════════╗
% ║           for figures via inkspace                  ║
% ╚═════════════════════════════════════════════════════╝

% to import figures via: https://github.com/gillescastel/inkscape-figures
\usepackage{import}
\usepackage{pdfpages}
\usepackage{transparent}


% This command is the ONLY command imported via the packages snippet, think of it as a tiny package written out
\newcommand{\incfig}[2][1]{%
    \def\svgwidth{#1\columnwidth}
    \import{./figures/}{#2.pdf_tex}
}

% ╔═══════════════════════════════════════════════════════════╗
% ║                 bibliography and citing                   ║
% ╚═══════════════════════════════════════════════════════════╝

\usepackage{csquotes}
% \usepackage[backend=biber,citestyle=alphabetic,bibstyle=authortitle]{biblatex}
\usepackage{biblatex}

\addbibresource{bibliography.bib}

% ╔═══════════════════════════════════════════════════════════╗
% ║                    colours (colors)                       ║
% ╚═══════════════════════════════════════════════════════════╝

\definecolor{dkgreen}{rgb}{0,0.6,0}
\definecolor{gray}{rgb}{0.5,0.5,0.5}
\definecolor{mauve}{rgb}{0.58,0,0.82}
\definecolor{lightyellow}{rgb}{1,1,0.6}
\definecolor{reddish}{HTML}{A01A3A}

%  ╔══════════════════════════════════════════════════════════╗
%  ║                       book format                        ║
%  ╚══════════════════════════════════════════════════════════╝

\pagestyle{fancy}

\setlength\headheight{20pt}

\renewcommand{\sectionmark}[1]{\markright{\thesection.\ #1}}

\fancyfoot[C,CO]{\textcolor{gray}{Nathanael Chwojko-Srawley}}
\fancyfoot[R,RO]{\thepage}{}

% Create a new page style for the first page
\fancypagestyle{firstpage}{
  \fancyhf{} % Clear header and footer
  \fancyhead[L]{\theauthor}
  \fancyhead[R]{\today} % Date in the top-right
  \fancyfoot[C]{\textcolor{gray}{Nathanael Chwojko-Srawley}}
  \fancyfoot[R]{\thepage}
}

% Apply the first page style conditionally
\AtBeginDocument{\thispagestyle{firstpage}}

%have format the title command
\newcommand{\chapter}[1]{%
  \clearpage
  \vspace*{30pt} % Space before the chapter title
  {\normalfont\huge\bfseries #1\par} % Chapter title formatting
  \vspace{20pt} % Space after the chapter title
}


\setlength{\parindent}{0pt} % don't like indentation infront of paragraphs
\setlength{\parskip}{1ex}   %so that there is still space between paragarphs


%  ╔══════════════════════════════════════════════════════════╗
%  ║                       Shortcuts                          ║
%  ╚══════════════════════════════════════════════════════════╝

%\ExerciseLevelInToc

% I like original mathcal better, but need lowercase.
\DeclareFontFamily{OT1}{pzc}{}
\DeclareFontShape{OT1}{pzc}{m}{it}{<-> s * [1.10] pzcmi7t}{}
\DeclareMathAlphabet{\mathpzc}{OT1}{pzc}{m}{it}


% mathbb
\newcommand{\A}{{\mathbb{A}}}
\newcommand{\C}{{\mathbb{C}}}
\newcommand{\F}{{\mathbb{F}}}
\newcommand{\N}{{\mathbb{N}}}
\newcommand{\Q}{{\mathbb{Q}}}
\newcommand{\R}{{\mathbb{R}}}
\newcommand{\V}{{\mathbb{V}}}
\newcommand{\Z}{{\mathbb{Z}}}
\newcommand{\BI}{{\mathbb{I}}}
\renewcommand{\H}{{\mathbb{H}}}
\renewcommand{\P}{{\mathbb{P}}}


% mathcal
% \newcommand{\co}{{\mathfrak{o}}} %mathcal{o} looks like wreath product, so I changed it to mathfrak
\newcommand{\co}{{\mathpzc{o}}}
\newcommand{\CA}{{\mathcal{A}}}
\newcommand{\CB}{{\mathcal{B}}}
\newcommand{\CC}{{\mathcal{C}}}
\newcommand{\CD}{{\mathcal{D}}}
\newcommand{\CF}{{\mathcal{F}}}
\newcommand{\CG}{{\mathcal{G}}}
\newcommand{\CH}{{\mathcal{H}}}
\newcommand{\CI}{{\mathcal{I}}}
\newcommand{\CK}{{\mathcal{K}}}
\newcommand{\CL}{{\mathcal{L}}}
\newcommand{\CM}{{\mathcal{M}}}
\newcommand{\CN}{{\mathcal{N}}}
\newcommand{\CO}{{\mathcal{O}}}
\newcommand{\CQ}{{\mathcal{Q}}}
\newcommand{\CR}{{\mathcal{R}}}
\newcommand{\CS}{{\mathcal{S}}}
\newcommand{\CT}{{\mathcal{T}}}
\newcommand{\CV}{{\mathcal{V}}}
\newcommand{\CZ}{{\mathcal{Z}}}
% EXCPETION:
\newcommand{\CP}{{\mathcal{P}}}
% \newcommand{\CP}{\mathbb{C}\mathrm{P}}
\newcommand{\RP}{\mathbb{R}\mathrm{P}}


%Euscript
\newcommand{\EB}{\EuScript{B}}
\newcommand{\EC}{{\EuScript{C}}}
\newcommand{\EF}{{\EuScript{F}}}
\newcommand{\EL}{{\EuScript{L}}}
\newcommand{\EO}{{\EuScript{O}}}


%mathfrak (lowercase)
\newcommand{\fa}{{\mathfrak{a}}}
\newcommand{\fb}{{\mathfrak{b}}}
\newcommand{\fc}{{\mathfrak{c}}}
\newcommand{\fd}{{\mathfrak{d}}}
\newcommand{\fm}{{\mathfrak{m}}}
\newcommand{\fn}{{\mathfrak{n}}}
\newcommand{\fo}{{\mathfrak{o}}}
\newcommand{\fp}{{\mathfrak{p}}}
\newcommand{\fq}{{\mathfrak{q}}}


%mathfrak (upper case)
\newcommand{\FA}{{\mathfrak{A}}}
\newcommand{\FB}{{\mathfrak{B}}}
\newcommand{\FC}{{\mathfrak{C}}}
\newcommand{\FD}{{\mathfrak{D}}}
\newcommand{\FK}{{\mathfrak{K}}}
\newcommand{\FM}{{\mathfrak{M}}}
\newcommand{\FN}{{\mathfrak{N}}}
\newcommand{\FO}{{\mathfrak{O}}}
\newcommand{\FP}{{\mathfrak{P}}}


%mathscr
\newcommand{\ff}{{\mathscr{F}}}
\newcommand{\SA}{\mathscr{A}}
\newcommand{\SC}{\mathscr{C}}
\newcommand{\SF}{\mathscr{F}}
\newcommand{\SG}{\mathscr{G}}
\newcommand{\SH}{\mathscr{H}}
\newcommand{\SI}{\mathscr{I}}
\newcommand{\SK}{\mathscr{K}}
\newcommand{\SN}{\mathscr{N}}
\newcommand{\SQ}{\mathscr{Q}}
\newcommand{\SR}{\mathscr{R}}
\newcommand{\SV}{\mathscr{V}}
\newcommand{\SZ}{\mathscr{Z}}


% operatorname -- 2 letters (lowercase and upper case)
\newcommand{\ab}{{\operatorname{ab}}}
\newcommand{\ad}{{\operatorname{ad}}}
\newcommand{\ch}{{\operatorname{ch}}}
\newcommand{\ev}{{\operatorname{ev}}}
\newcommand{\id}{{\operatorname{id}}}
\newcommand{\im}{{\operatorname{im}}}
\newcommand{\nm}{{\operatorname{Nm}}}
\newcommand{\op}{{\operatorname{op}}}
\newcommand{\rk}{{\operatorname{rk}}}
\newcommand{\sh}{{\operatorname{sh}}}
\newcommand{\tr}{{\operatorname{tr}}}

\newcommand{\Ab}{{\operatorname{Ab}}}
\newcommand{\Br}{{\operatorname{Br}}}
\newcommand{\Ch}{{\operatorname{Ch}}}
\newcommand{\Cl}{{\operatorname{Cl}}}
\newcommand{\GL}{{\operatorname{GL}}}
\newcommand{\Nm}{{\operatorname{Nm}}}
\newcommand{\SL}{{\operatorname{SL}}}

\renewcommand{\Re}{{\operatorname{Re}}}
\renewcommand{\Im}{{\operatorname{Im}}}


%categories
\newcommand{\Grp}{{\textbf{Grp}}}
\newcommand{\Fld}{{\textbf{Fld}}}
\newcommand{\Sh}{{\textbf{Sh}}}
\newcommand{\Set}{{\textbf{Set}}}
\newcommand{\Mod}{{\textbf{Mod}}}
\newcommand{\PSh}{{\textbf{PSh}}}
\newcommand{\Sch}{{\textbf{Sch}}}
\newcommand{\Top}{{\textbf{Top}}}
\newcommand{\RgSp}{{\textbf{RgSp}}}
\newcommand{\Ring}{{\textbf{Ring}}}


%operatorname -- 3+ letters (lowercase and upper case)
\newcommand{\rad}{{\operatorname{rad}}}
\newcommand{\sep}{{\operatorname{sep}}}
\newcommand{\res}{{\operatorname{res}}}
\newcommand{\sgn}{{\operatorname{sgn}}}
\newcommand{\pre}{{\operatorname{pre}}}
\newcommand{\ord}{{\operatorname{ord}}}
\newcommand{\vol}{{\operatorname{vol}}}
\newcommand{\lcm}{{\operatorname{lcm}}}

\newcommand{\conj}{{\operatorname{conj}}}
\newcommand{\disc}{{\operatorname{disc}}}
\newcommand{\stab}{{\operatorname{stab}}}
\newcommand{\kdim}{{\operatorname{kdim}}}
\newcommand{\diag}{{\operatorname{diag}}}

\newcommand{\trdeg}{\operatorname{trdeg}}
\newcommand{\coker}{{\operatorname{coker}}}
\newcommand{\codim}{{\operatorname{codim}}}

\newcommand{\Ann}{{\operatorname{Ann}}}
\newcommand{\Aut}{{\operatorname{Aut}}}
\newcommand{\Der}{{\operatorname{Der}}}
\newcommand{\Div}{{\operatorname{Div}}}
\newcommand{\Emb}{{\operatorname{Emb}}}
\newcommand{\End}{{\operatorname{End}}}
\newcommand{\Ext}{{\operatorname{Ext}}}
\newcommand{\Gal}{{\operatorname{Gal}}}
\newcommand{\Hom}{{\operatorname{Hom}}}
\newcommand{\Ind}{{\operatorname{Ind}}}
\newcommand{\Inn}{{\operatorname{Inn}}}
\newcommand{\Int}{{\operatorname{int}}}
\newcommand{\Irr}{{\operatorname{Irr}}}
\newcommand{\Jac}{{\operatorname{Jac}}}
\newcommand{\Mor}[1]{\operatorname{Mor}(\textbf{#1})}
\newcommand{\Nil}{{\operatorname{Nil}}}
\newcommand{\Obj}{{\operatorname{Obj}}}
\newcommand{\Out}{{\operatorname{Out}}}
\newcommand{\Pic}{{\operatorname{Pic}\,}}
\newcommand{\Rep}{{\operatorname{Rep}}}
\newcommand{\Res}{{\operatorname{Res}}}
\newcommand{\Spl}{{\operatorname{Spl}}}
\newcommand{\Syl}{{\operatorname{Syl}}}
\newcommand{\Sym}{{\operatorname{Sym}}}
\newcommand{\Tor}{{\operatorname{Tor}}}

\newcommand{\Char}{{\operatorname{char}}}
\newcommand{\Frac}{{\operatorname{Frac}}}
\newcommand{\Frob}{{\operatorname{Frob}}}
\newcommand{\Proj}{{\operatorname{Proj}\,}}
\newcommand{\RMod}{{}_R\operatorname{Mod}}
\newcommand{\SExt}{{\mathcal{E}\emph{xt}}}
\newcommand{\SHom}{{\emph{Hom}}}
\newcommand{\Span}{{\operatorname{span}}}
\newcommand{\Spec}{{\operatorname{Spec}\,}}
\newcommand{\Supp}{{\operatorname{Supp}}}
\newcommand{\Tens}[1]{#1\otimes --}
\newcommand{\fHom}[1]{\operatorname{Hom}(#1, --)}

\newcommand{\mSpec}{{\operatorname{mSpec}}}
\newcommand{\cofHom}[1]{\operatorname{Hom}(--, #1)}


% connectors
% \renewcommand{\mid}{\big|}
\newcommand{\sse}{\subseteq}
\newcommand{\subg}{\leqslant}
\newcommand{\supg}{\geqslant}
\newcommand{\ssne}{\subsetneq}
\newcommand{\isosubg}{\lesssim}
\newcommand{\nsse}{\not\subseteq}
\newcommand{\nsubg}{\trianglelefteq}
\newcommand{\nsupg}{\trianglerighteq}
\newcommand{\csubg}{\blacktriangleleft}
\renewcommand{\mid}{\big|}


% Arrows
\newcommand{\rw}{\rightarrow}
\newcommand{\Rw}{\Rightarrow}
\newcommand{\lw}{\leftarrow}
\newcommand{\Lw}{\Leftarrow}
\newcommand{\lrw}{\leftrightarrow}
\newcommand{\LRw}{\Leftrightarrow}
\newcommand{\acton}{\curvearrowright}
\newcommand{\actson}{\curvearrowright}
\newcommand\mapsfrom{\mathrel{\reflectbox{\ensuremath{\mapsto}}}}

% arrows for commutative diagram: create four new angled double-struck arrows
\newcommand\myrot[1]{\mathrel{\rotatebox[origin=c]{#1}{$\Rightarrow$}}}
\newcommand\NEarrow{\myrot{45}}
\newcommand\NWarrow{\myrot{135}}
\newcommand\SWarrow{\myrot{-135}}
\newcommand\SEarrow{\myrot{-45}}


%shorthands
\newcommand{\oln}[1]{{\overline{#1}}}
\renewcommand{\bf}[1]{\textbf{#1}}
\newcommand{\qn}{\quad\newline}


% limits
\newcommand{\Lim}{{\varprojlim}}
\newcommand{\coLim}{{\varinjlim}}
\newcommand{\Dlim}{\lim\limits_{\longrightarrow}}
\newcommand{\coDlim}{\lim\limits_{\longleftarrow}}


%for saturated ideals in the localization section (I don't like these, I think I'll delete)
\newcommand{\LIm}{{}^e }
\newcommand{\LPre}{{}^c }

%  ╔══════════════════════════════════════════════════════════╗
%  ║                    Custom Commands                       ║
%  ╚══════════════════════════════════════════════════════════╝

% swap varphi and phi
\let\temp\phi
\let\phi\varphi
\let\varphi\temp

\newcommand{\placeholder}{\_\_\_\_\_}

\newcommand{\set}[2]{\left\{#1 \ \middle|\ #2\right\}}

%mod should not have the crazy amount of space to its right!
\renewcommand{\mod}{\,\operatorname{mod}\,}

% dcups
\newcommand{\dcup}{\sqcup}
\newcommand{\Dcup}{\coprod}
% \newcommand{\Dcup}{\bigsqcup}

\newcommand{\nbot}{\mathrel{\rlap{$\bot$}\mkern2mu/}}

%a nice large circle
\newcommand{\tikzcircle}[2][red,fill=black]{\tikz[baseline=-0.5ex]\draw[#1,radius=#2] (0,0) circle ;}%

\makeatletter
\newcommand*\bigcdot{\mathpalette\bigcdot@{.5}}
\newcommand*\bigcdot@[2]{\mathbin{\vcenter{\hbox{\scalebox{#2}{$\m@th#1\bullet$}}}}}
\makeatother

%somehow not a default command
\def\rddots{\mathstrut^{.^{.^{.}}}}

%% new delimiter
\DeclarePairedDelimiter{\ceil}{\lceil}{\rceil}


%  ╔══════════════════════════════════════════════════════════╗
%  ║                     environments                         ║
%  ╚══════════════════════════════════════════════════════════╝


%remember the option: breakable

% create tcolor boxes
\newtcbtheorem[number within=section]{thm}{Theorem}%
{
colback=green!5,
colframe=green!35!black,
fonttitle=\bfseries,
label type=theorem
}{th}

\newtcbtheorem[number within=section, use counter from=thm]{lem}{Lemma}%
{
colback=blue!5,
colframe=black!35!black,
fonttitle=\bfseries,
label type=lemma
}{lm}
\newtcbtheorem[number within=section, use counter from=thm]{prop}{Proposition}%
{
colback=white!5,
colframe=white!35!black,
fonttitle=\bfseries,
label type=proposition
}{pr}
\newtcbtheorem[number within=section, use counter from=thm]{cor}{Corollary}%
{colback=blue!5,
colframe=blue!35!black,
fonttitle=\bfseries,
label type=corollary
}{co}
\newtcbtheorem[number within=section, use counter from=thm]{axiom}{Axiom}%
{colback=black!5,
colframe=yellow!35!black,
fonttitle=\bfseries,
label type=axiom
}{ax}
\newtcbtheorem[number within=section, use counter from=thm]{defn}{Definition}%
{colback=red!5,
colframe=red!35!black,
fonttitle=\bfseries,
label type=definition
}{df}


% for <s-cr> to create \item[$\LRw$] instead of \item
\newenvironment{equivEnumerate}
 {\begin{enumerate}
	}{
\end{enumerate}}



%Custom enviroments
 \newenvironment{note}
 {\begin{tcolorbox}[enhanced, sharp corners, colback=white , borderline={0.2pt}{0pt}{black}]
   \begin{quote}\textbf{Note}:
   }{
   \end{quote}\end{tcolorbox}}

 \newenvironment{intuition}
 {\begin{quote}\textbf{Intuition}:
   }{
   \end{quote}}

\newtcbtheorem[number within=section, use counter from=thm]{titledBox}{}{%
  enhanced,
  sharp corners,
  colback=white,
  colframe=black,
  borderline={0.2pt}{0pt}{black},
  left=2mm,
  right=2mm,
  top=2mm,
  bottom=2mm,
  title style={opacity=1},
  attach title to upper,
  before upper={\textbf{\tcbtitletext}\quad},
  coltitle=black,
  fonttitle=\bfseries,
  separator sign={\quad}
}{box}


%better example and proof environments
\def\exampletext{Example} % If English
\colorlet{colexam}{red!55!black} % Global example color


\newtcbtheorem[number within=section, use counter from=thm]{example}{Example}{% \newtcolorbox[use counter=testexample]{testexamplebox}{%
% Example Frame Start
empty,% Empty previously set parameters
title={\exampletext \thetcbcounter #1},% use \thetcbcounter to access the testexample counter text
% Attaching a box requires an overlay
attach boxed title to top left,
   % Ensures proper line breaking in longer titles
   minipage boxed title,
% (boxed title style requires an overlay)
boxed title style={empty,size=minimal,toprule=0pt,top=4pt,left=3mm,overlay={}},
coltitle=colexam,fonttitle=\bfseries,
before=\par\medskip\noindent,parbox=false,boxsep=0pt,left=3mm,right=0mm,top=2pt,breakable,pad at break=0mm,
   before upper=\csname @totalleftmargin\endcsname0pt, % Use instead of parbox=true. This ensures parskip is inherited by box.
% Handles box when it exists on one page only
overlay unbroken={\draw[colexam,line width=.5pt] ([xshift=-0pt]title.north west) -- ([xshift=-0pt]frame.south west); },
% Handles multipage box: first page
overlay first={\draw[colexam,line width=.5pt] ([xshift=-0pt]title.north west) -- ([xshift=-0pt]frame.south west); },
% Handles multipage box: middle page
overlay middle={\draw[colexam,line width=.5pt] ([xshift=-0pt]frame.north west) -- ([xshift=-0pt]frame.south west); },
% Handles multipage box: last page
overlay last={\draw[colexam,line width=.5pt] ([xshift=-0pt]frame.north west) -- ([xshift=-0pt]frame.south west); }%
}{ex}



\def\prooftext{Proof} % If English
\NewDocumentEnvironment{Proof}{ O{} }
{
    \colorlet{colexam}{black!55!black} % Global example color
    \newtcolorbox[number within=section]{testproofbox}{%
        empty,% Empty previously set parameters
        title={\emph{\prooftext} \thetcbcounter: #1},
        attach boxed title to top left,
        minipage boxed title,
        boxed title style={empty,size=minimal,toprule=0pt,top=4pt,left=3mm,overlay={}},
        coltitle=colexam,
        fonttitle=\bfseries,
        before=\par\medskip\noindent,
        parbox=false,
        boxsep=0pt,
        left=3mm,
        right=0mm,
        top=2pt,
        breakable,
        pad at break=0mm,
        before upper=\csname @totalleftmargin\endcsname0pt,
        % Removed height constraints
        % Added flexible width
        width=\dimexpr\textwidth-3mm\relax,
        % Modified overlays
        overlay unbroken={\draw[colexam,line width=.5pt] ([xshift=-0pt]title.north west) -- ([xshift=-0pt]frame.south west); },
        overlay first={\draw[colexam,line width=.5pt] ([xshift=-0pt]title.north west) -- ([xshift=-0pt]frame.south west); },
        overlay middle={\draw[colexam,line width=.5pt] ([xshift=-0pt]frame.north west) -- ([xshift=-0pt]frame.south west); },
        overlay last={\draw[colexam,line width=.5pt] ([xshift=-0pt]frame.north west) -- ([xshift=-0pt]frame.south west); },%
    }
    \begin{testproofbox}
}
{\end{testproofbox}\endlist}


%  ╔══════════════════════════════════════════════════════════╗
%  ║                    for Dynkin diagrams                   ║
%  ╚══════════════════════════════════════════════════════════╝

\def\row#1/#2!{#1_{\IfStrEq{#2}{}{n}{#2}} & \dynkin{#1}{#2}\\}
\newcommand{\tble}[1]{
   \renewcommand*\do[1]{\row##1!}
   \[
      \begin{array}{ll}\docsvlist{#1}\end{array}
   \]
}

%  ╔══════════════════════════════════════════════════════════╗
%  ║                         links                            ║
%  ╚══════════════════════════════════════════════════════════╝

\definecolor{LinkColour}{HTML}{A64A3B} % favourite
% \definecolor{LinkColour}{HTML}{8C3D32} % 2nd place
\hypersetup{
  colorlinks,
  citecolor=black,
  filecolor=magenta,
  linkcolor=LinkColour,
  urlcolor=blue,
}


%  ╔══════════════════════════════════════════════════════════╗
%  ║                          misc                            ║
%  ╚══════════════════════════════════════════════════════════╝

\stackMath %to be able to stack text in math mode
\makeindex

%import options: all, bibliography, book-formatting, booksSetting, customEnvironments, dynkin, hw-formatting, language-setup, newcommands, packages, preamble, proofs, settings, tcolorbox, test.svg,
\begin{document}
\maketitle


This blog is an extract from EYTNKA Algebra \cite{NateAlgebra} going over my intuition on why field extensions
is the language of galois theory

\qn
For eons, mathematicians where looking for relations between coefficients and roots of a polynomial. Some of
these were found with success, for example if $x^2+bx+c$ has roots $\alpha$ and $\beta$, then $\alpha + \beta
= -b$ and $\alpha\beta = c$. More generally, a lot of focus was put into find the roots of such polynomials,
which lead to the discovery of the quadratic equation\footnote{This is a little anachronistic, for when the
  quadratic formula was first discovered in Babylonia and more succinctly by Brahmagupta, the ``modern''
  notion of polynomial and variables where not used, but rather it was shown by relating areas of circles and
rectangles}, and eventually the cubic and quartic equation. However, the effort to understand these brought
limited results to pressing questions like the solvability of the quintic, that is whether its roots can
written in terms of it's coefficients and the arithmetic operations of $+, -, \times, \div$ and $\sqrt{-}$. It
wasn't until the turn of the 18th century when serious progress has been made on this front. The key insight
was to consider breaking down a polynomials as though it already had roots into the product of monomials and
analyzing the relationship. The first immediate realization that could be made linking roots to  coefficients
is that given we can completely split a polynomial:
\[
	x^n + \cdots + \alpha_1x + \alpha_0 = (x-a_1)(x-a_2)\cdots(x-a_n)
\]
each coefficient of the original polynomial can be represented as a \emph{symmetric function}:
\begin{align*}
	\alpha_{n-1} &= a_1 + a_2 + \cdots + a_n\\
	\alpha_{n-2} &= a_1a_2 + a_1a_3 + \cdots + a_1a_n + a_2a_3 + \cdots a_{n-1}a_n\\
		 &\cdots\\
	\alpha_0 &= a_1a_2\cdots a_n
\end{align*}
These form a natural relationship between the roots and the coefficients. It was noticed that these equation
are invariant under permutation action of $S_n$; if $A = \{A_0, ..., A_{n-1}\}$ represent the above $n$
equations and the action is given by $(r\ s)\cdot a_r = a_s$, then $\sigma \cdot A_i = A_i$ for each $A_i$ and
$\sigma$, hence the action $S_n \actson A$ is the trivial action. Thus, there is always a natural
group-action.

As any roots of polynomial induce such equations and induce this trivial action, these don't give enough
information about polynomials. It shall be shown that roots of a generic polynomial satisfy no further
equation, and in these cases there is nothing more that can be done to understand a polynomial via a
group-action on its roots. However, there are polynomials whose roots shall have \emph{more} relations the
roots satisfy, limiting what group can act on them (i.e. a subgroup of $S_n$ will act trivially, but $S_n$
will not). This chapter shall make rigorous this realization and translate it into the modern language of
field theory.

This realization almost managed to show the general quintic is unsolvable:

\begin{quote}
	``The quintic was almost proven to have no general solutions by radicals by Paolo Ruffini in 1799,
	whose key insight was to use permutation groups, not just a single permutation. His solution
	contained a gap, which Cauchy considered minor, though this was not patched until the work of the
	Norwegian mathematician Niels Henrik Abel, who published a proof in 1824, thus establishing the
	Abel–Ruffini theorem.''\footnote{taken from Wikipedia for Galois Theory}
\end{quote}

This theory was established with the aim to show there is an unsolvable  quintic, but not how to find whether a given quintic is
solvable. This is where the work of \'Evariste Galois comes in. He noticed that the key insight was to work
with all possible relationships of the roots of the polynomial, assuming all the roots exist.


\subsection*{From action on roots to field theory}

Take the polynomial:
\[
	x^4-2=0
\]
which is irreducible over $\Q$ by Eisenstein's criterion. Adjoining
$\sqrt[4]{2}$, $\Q(\sqrt[4]{2})$, the polynomial splits to:
\[
	(x-\sqrt[4]{2})(x+\sqrt[4]{2})(x^2+\sqrt{2})
\]
Adjoining $\sqrt[4]{2}i$, the polynomial splits completely:
\begin{equation}\label{eq:forthRootOf2}
	(x-\sqrt[4]{2})(x+\sqrt[4]{2})(x-\sqrt[4]{2}i)(x+\sqrt[4]{2}i)
\end{equation}
Already, there is some interesting phenomena at play, for example, another way of getting this splitting is by adjoining
$\sqrt[4]{2}$ and then adjoining $i$ (or adjoining $i$ then $\sqrt[4]{2}$). In
fact, by just adjoining $\sqrt[4]{2} + i$, we immediately get this splitting. Indeed, certainly $\Q(\sqrt[4]{2} + i)
\sse \Q(\sqrt[4]{2}, i)$. Conversely, to show $i \in \Q\sqrt[4]{2}+i$, letting $a = \sqrt[4]{2}+i$, then:
\begin{align*}
	a-i &= \sqrt[4]{2}\\
	\LRw\ (a-i)^4 &= 2\\
	\LRw\ (a-i)^4 &= 2\\
	\LRw\ a^4 - 4a^3i - 6a^2 + 4ai + 1 &= 2\\
	\LRw\  i(4a- 4a^3i) &= 1-a^4+ 6a^2\\
	\LRw\  i &= \frac{1-a^4+ 6a^2}{4a- 4a^3i}\\
\end{align*}
showing that $i \in \Q(\sqrt[4]{2}+i)$, and so $a-i = \sqrt[4]{2} \in \Q(\sqrt[4]{2}+i)$. More generally for
any finite field separable extension $K/\Q$, such an element always exists and is constructible by the
primitive element theorem.

Let us understand the equations the roots form given by symmetric functions. As $x^4-2 = x^4 + 0x^3 + 0x^2 + 0x-2$ The only
meaningful relation obtained by multiplying out the roots in equation~(\ref{eq:forthRootOf2}) is:
\[
	(\sqrt[4]{2})(-\sqrt[4]{2})(\sqrt[4]{2}i)(-\sqrt[4]{2}i) = 2 \qquad \text{equiv.} \qquad
	(\sqrt[4]{2})(-\sqrt[4]{2})(\sqrt[4]{2}i)(-\sqrt[4]{2}i)-2 = 0
\]
This equation is invariant of the action of $S_4$, namely if
\[
	(a_1, a_2, a_3, a_4)  = (\sqrt[4]{2}, \sqrt[4]{2}i, -\sqrt[4]{2}, -\sqrt[4]{2}i)
\]
represent the $4$ roots, and $A_0 = (\sqrt[4]{2})(-\sqrt[4]{2})(\sqrt[4]{2}i)(-\sqrt[4]{2}i)-2 = 0$ , then:
\[
	\sigma\cdot A_0 = A_0 \qquad \forall \sigma \in S_4
\]
However, there are more equations that can be formed with these roots (using the operations of a field):
\begin{align*}
	A_1&:\ \sqrt[4]{2} - \sqrt[4]{2}  =0\\
	A_2&:\ \sqrt[4]{2}i - \sqrt[4]{2}i = 0\\
\end{align*}
These equations are \emph{not} invariant under the permutation action of $S_4$:
\[
	(1\ 2)\cdot A_1 \ne A_1 \qquad \text{(concretely: $(1\ 2) \cdot (\sqrt[4]{2} - \sqrt[4]{2})
	= (\sqrt[4]{2}i - \sqrt[4]{2}) \ne A_1$)}
\]
A subgroup of $S_4$ must be chosen, namely notice that $D_4 \cong \langle (1\ 2\ 3\ 4), (1\ 3) \rangle \sse
S_4$ keeps these equations invariant:
\[
	D_4\cdot A_0 = A_0\qquad D_4 \cdot A_1 = A_1 \qquad D_4 \cdot A_2 = A_2
\]
Are there any other equations that we can be considered which would further restrict the choice of invariant
permutation action? To answer this, the above is re-framed in the language of field extensions.

Let us say $f$ is a polynomial over $\Q$, $\alpha$ is
a root. Then the algebraic extension $\Q(\alpha)/\Q$ contains \emph{all} possible polynomial equations that can be
generated with elements from $\Q$ and $\alpha$ (recall $\Q(\alpha) = \Q[\alpha]$). If any of these polynomials happen to be another
root of $f$, that is $p(\alpha) = \beta$ where $f(\beta) = 0$ and $p(x) \in \Q[x]$, then $\Q(\alpha)/\Q$ contains
this relationship between the roots. Let's say one of these relationships is represented as :
\[
	P(x) = p(x) - \beta \text{ so that } P(\alpha) = p(\alpha)-\beta = 0
\]
Then a $\Q$-automorphism (in this case nothing more than a field automorphism) would have to preserve this
relationship (it must map roots to roots). Hence, the set of $\Q$-automorphism captures the relationships
between $\alpha$ and all possible roots that are $\Q$-polynomials of $\alpha$. If $K/\Q$ is the splitting
field of $K$, then the $\Q$-automorphisms of $K$ would be restricted to automorphism that must satisfy all
possible relationships between roots. Given $K/\Q$ is a splitting field for $f$, as any $\Q$-automorphisms
$\sigma$ is determined by where the roots map to, it is associated to a permutation in $S_n$. If there were no
relationships (you need to add $n$ roots to split $f$), then the automorphism group would be isomorphic to
$S_n$. If there is any other type of group, that means there are some relationships between the coefficients
(just like the example of $\Q(\sqrt[4]{2} + i)/\Q$).

The automorphism group of the splitting field above is called the \emph{galois group}, and
in~\cite{NateAlgebra} it is used to answer some long-pending mathematical questions such as the
solvability of [general] quintic.

Another way of thinking about galois group is they are a rough classification tool for field
extensions: if two field extensions have different Galois groups they cannot be the same extensions. It is
similar in roughness to the fundamental group $\pi_1(X)$ as is often encountered in algebraic topology.
To see this roughness in classification, as the splitting field
$\Q(\sqrt[4]{2}, i)$ of $x^4-2$ is equal to $\Q(\sqrt[4]{2} + i)$, and the minimal polynomial of
$\sqrt[4]{2}+i$ is:
\[
	x^8 + 4 x^6 + 2 x^4 + 28 x^2 + 1
\]
then they share the same restrictions on their roots as they have the same galois group, but they certainly
also have many differences (for example, these two polynomials have different discriminants).

\begin{titledBox*}{Galois Group as Fundamental Group}
	\qn
	(Assuming strong familiarity of algebraic geometry and algebraic topology) Under the right framework in algebraic geometry,
	galois groups can be interpreted as \emph{\'etale fundamental group} for $\Spec K$,
	\[
		\pi^{\text{\'et}}_1(\Spec(K)) \cong \Gal_F(\overline{F})
		\]
	The name \emph{fundamental group} is well-deserved for if $X$
	is a connected $\C$-scheme of finite type (i.e. ``essentially'' a connected variety) and $X(\C)$ is it's
	analytification, then:
	\[
		\pi^{\text{\'et}}_1(X) \cong \widehat{\pi_1^{\operatorname{top}}(X(\C))}
	\]
	where $\pi_1^{\operatorname{top}}(-)$ is the usual fundamental group and $\widehat{-}$ is the profinite
	completion with respect to its normal subgroups. See~\cite{NateAlgGeo} for the details.
\end{titledBox*}


\subsection*{Modern uses of Galois Theory}

Galois groups are central in both algebraic geometry and number theory, where they serve different purposes in
each domain. The galois group of a polynomial:
\begin{enumerate}
	\item establishes the algebraic relations between the roots and creates a ``covering space'' that has some
		similar properties to the covering spaces seen in algebraic topology (as remarked in the box above),
	\item establishes the invariance necessary to find a solution over a lower field, in particular if $L/K$
		is a field extension and a root of an $L[x]$ polynomial has a root in $K$, then the root is invariant under the
		$\Gal_K(L)$-action.
\end{enumerate}
\newpage
% \newpage
\printbibliography

\end{document}
